%% Conclusion section

\section{Conclusion}
\label{sec:conclusion}
We presented \approach{}, a retrieval-augmented generation (RAG) framework for DeFi attack detection that addresses key limitations of existing approaches. \approach{} constructs an incrementally expandable knowledge base of historical attacks with behavioral summaries, employs a two-stage retrieval mechanism with dynamic similarity thresholding, \Blue{applies adaptive prompting strategies, and incorporates XGBoost post-filtering.} Together, these components bridge the gap between low-level execution traces and high-level attack semantics, while mitigating LLM hallucinations via evidence-grounded reasoning.

\Blue{\approach{} achieves 92.93\% F1 on \dsystem{}, our 728-transaction evaluation set constructed from DeFiHackLabs attacks and sampled benign transactions and reduces false positives from 60 to 7 compared with XGBoost-only, an F1 improvement of 6.45 percentage points. The complete system also exceeds the no-XGBoost variant by 14.99 percentage points in F1. \textcolor{red}{Without XGBoost, two-stage retrieval increases precision but reduces recall and F1.} The complementary evaluation on the 91-attack subset of DeFiScope\textquotesingle s D1 dataset~\cite{zhong2025detectingvariousdefiprice} detects 86 attacks. The combined pipeline uses historical attack evidence to support LLM analysis and numerical features to check its attack candidates. Its advantage over direct numerical classification is a substantial reduction in false alarms at comparable recall.}

\approach{} faces limitations in scenarios involving cryptographic exploits without behavioral signatures, complex multi-step cross-contract attacks with fragmented semantic signals, and truly novel attack patterns that are absent from the knowledge base. Future work includes integrating formal verification for cryptographic exploit detection, developing more sophisticated call-graph analysis for cross-contract reasoning, and implementing automated knowledge base expansion to remain responsive to emerging threats.

\approach{} demonstrates that retrieval-augmented generation is a promising paradigm for security analysis tasks that require both domain knowledge and robust reasoning. As DeFi protocols continue to evolve, approaches that leverage accumulated security knowledge while adapting to new patterns will be essential for protecting the assets at stake in this ecosystem.
